\scanpage{68}
\mspage{C}
% archive foliation: 67
Reprenons notre $Z$ normal de dim $n\geq2$, et $X=Z-Y$.
Plongeons $Z$ dans $P^r$, et le morphisme induit $f:X\to P^r$.
On prend $s=n-1$, de sorte que
$X_u$ \deleted{est une \uline{courbe}}
\inserted{$=X_K\cap L_u=Z_K\cap L_u$, car $Y_K\cap L_u=\varnothing$,}
donc $X_u=Z_K\cap L_u$ est une \uline{courbe} de
\deleted{genre $g$} projective, d'ailleurs normale.
Si $g$ est son genre, il résulte alors du corollaire au lemme de Bertini
que $\pi_1(X)$ admet $2g$ générateurs topologiques, OK.

\uline{Conséquence.} Soit $K$ une extension de type fini d'un corps $k$.
Il existe tjrs des modèles projectifs,
\deleted{\unsure{(Alors normaux}}
à fortiori propres, et normaux si on veut,
\deleted{donc réguliers en codim $\leq1$,}
\deleted{\inserted{Soit $X$ \unsure{le}}
\unsure{les modèles de tels}}
\deleted{modèles réguliers $U$ qui sont
\unsure{gros}
``\uline{presque propres}'', i.e. complémentaires
d'un ens fermé de codim $\geq2$ dans un modèle propre $X$.
D'ailleurs, tts les modèles $U$ comme dessus, ordonnés par}

\margin{\diagram{\begin{tikzcd}
&X''\arrow[dl,no head]\arrow[dr,no head]&\\
X&&X'
\end{tikzcd}}}
